\documentclass[a4paper]{article}

%% Language and font encodings
\usepackage[spanish]{babel}
\usepackage[utf8x]{inputenc}
\usepackage[T1]{fontenc}

%% Sets page size and margins
\usepackage[a4paper,top=3cm,bottom=2cm,left=2cm,right=2cm,marginparwidth=1.75cm]{geometry}

%% Useful packages
\usepackage{amsmath}
\usepackage{graphicx}
\usepackage{caption}
\usepackage{subcaption}
\usepackage{float}
\usepackage{titling} 
\usepackage{booktabs} % Tablas
\usepackage{pbox} % Dos líneas en tabla
%\usepackage{subfigure} % Imágenes
%\usepackage{subcaption} % Pie de imagen

%% Referencias/Bibliografía
%\usepackage[backend=biber, style=apa]{biblatex}
%\addbibresource{ref.bib} % Nombre del archivo con las referencias

%Links de referencias
\usepackage[hidelinks]{hyperref} % Hace que las referencias sean clicables
\newcommand{\myref}[1]{\textsuperscript{\textbf{\ref{#1}}}} % Define un estilo propio

% Código
\usepackage{listings}
\usepackage{xcolor}

% Definición de colores atractivos
\definecolor{codegreen}{rgb}{0,0.6,0}
\definecolor{codegray}{rgb}{0.5,0.5,0.5}
\definecolor{codepurple}{rgb}{0.58,0,0.82}
\definecolor{backcolour}{rgb}{0.95,0.95,0.92}
\definecolor{terminalback}{rgb}{0.1,0.1,0.1}
\definecolor{terminaltext}{rgb}{0.9,0.9,0.9}

% Configuración general para Python
\lstdefinestyle{pythonstyle}{
    backgroundcolor=\color{backcolour},   
    commentstyle=\color{codegreen},
    keywordstyle=\color{magenta},
    numberstyle=\tiny\color{codegray},
    stringstyle=\color{codepurple},
    basicstyle=\ttfamily\small,
    breakatwhitespace=false,         
    breaklines=true,                 
    captionpos=b,                    
    keepspaces=true,                 
    numbers=left,                    
    numbersep=5pt,                  
    showspaces=false,                
    showstringspaces=false,
    showtabs=false,                  
    tabsize=4,
    language=Python,
    literate= {á}{{\'a}}1 {é}{{\'e}}1 {í}{{\'i}}1 {ó}{{\'o}}1 {ú}{{\'u}}1 {ñ}{{\~n}}1 {Á}{{\'A}}1 {É}{{\'E}}1 {Í}{{\'I}}1 {Ó}{{\'O}}1 {Ú}{{\'U}}1 {Ñ}{{\~N}}1 % Para las tildes y caracteres del español
}

% Configuración para comandos de Terminal
\lstdefinestyle{terminalstyle}{
    backgroundcolor=\color{terminalback},   
    basicstyle=\ttfamily\small\color{terminaltext},
    breakatwhitespace=false,         
    breaklines=true,                 
    captionpos=b,                    
    keepspaces=true,                 
    numbers=none, % Sin números de línea en la terminal
    showspaces=false,                
    showstringspaces=false,
    showtabs=false,                  
    tabsize=4,
    literate= {á}{{\'a}}1 {é}{{\'e}}1 {í}{{\'i}}1 {ó}{{\'o}}1 {ú}{{\'u}}1 {ñ}{{\~n}}1 {Á}{{\'A}}1 {É}{{\'E}}1 {Í}{{\'I}}1 {Ó}{{\'O}}1 {Ú}{{\'U}}1 {Ñ}{{\~N}}1 % Para las tildes y caracteres del español
}

%%%%%%%% DOCUMENT %%%%%%%%
\begin{document}

%%%% Title Page
\begin{titlepage}

\newcommand{\HRule}{\rule{\linewidth}{0.5mm}} 							% horizontal line and its thickness
\center 
 
% University
\textsc{\LARGE Universidad de Granada}\\[1cm]

% Document info
\textsc{\Large Periféricos y Dispositivos de Interfaz Humana}\\[0.2cm]
\textsc{\large Curso 2025/2026}\\[1cm] 										% Course Code
\HRule \\[0.8cm]
{ \huge \bfseries Práctica 6. Reconocimiento de imágenes}\\[0.7cm]								% Assignment
\HRule \\[2cm]
\large
\emph{Autora:}\\
Inés Ruiz Sánchez\\[1.5cm]													% Author info
{\large \today}\\[5cm]

\includegraphics[width=0.3\textwidth]{images/UGRlogo.png}\\[1cm] 	% University logo
\vfill 
\end{titlepage}

% Generar índice de contenidos y figuras
\tableofcontents
\listoffigures
\lstlistoflistings

%%%% SECCIONES

\newpage
%% Sección I
\section{Introducción y objetivos}
Los \textbf{objetivos} concretos de esta práctica son: 
\begin{itemize}
    \item Conocer cómo funcionan las técnicas de reconocimiento de imágenes. 
    \item Hacer un programa sencillo para realizar detección facial. 
\end{itemize}

\textbf{Repositorio} en GitHub: \url{https://github.com/ruiz314/PDIH/tree/main/P6}.\\ 

%----------------------------------------------------------------
%\newpage
%% Sección II
\section{Cuestiones a resolver}
Para superar esta práctica se proponen los siguientes ejercicios: 
\begin{enumerate}
    \item Hacer reconocimiento de caras de personas en imágenes (ejemplo básico).
    \item Hacer reconocimiento de caras de personas en vídeos.
    \item Hacer reconocimiento de caras de gatos en vídeos.
    \item Hacer simultáneamente el reconocimiento de cuerpos y caras de personas en vídeos. Es requisito que ambas cosas se detenten al mismo tiempo. Como sugerencia, se puede hacer una primera aproximación para hacer la detección en imágenes y posteriormente en vídeos. Para este ejercicio se facilita un vídeo de ejemplo, aunque se pueden usar otros diferentes. 
\end{enumerate}


%----------------------------------------------------------------
\newpage
%% Sección III
\section{Configuración del entorno de trabajo}\label{sec:Configuracion}
% Configuración del entorno de trabajo
El sistema operativo usado es Microsoft Windows. A continuación, se detallan los pasos seguidos para la instalación del software, la gestión de entornos virtuales y la adquisición de los modelos de detección necesarios.


Tenía instalada la \textbf{versión 3.13.2 de Python}, como se puede ver en la Figura \ref{fig:python}.
\begin{figure}[h!]
    \centering
    \includegraphics[width=0.8\textwidth]{images/1_2python.png}
    \caption{Proceso de instalación de Python en Windows.}
    \label{fig:python}
\end{figure}

%Se ha creado un \textbf{entorno virtual} específico denominado \texttt{venv} dentro de la carpeta raíz del proyecto (\texttt{P6}) con el comando \texttt{python -m venv venv}. Esto es una práctica muy común cuando se trabaja con python. Una vez ejecutado, el prompt de la consola antepone el prefijo \texttt{(venv)}.\\

Para la instalación de la Librería \textbf{OpenCV-Python} uso la siguiente orden:

\begin{lstlisting}[style=terminalstyle, caption={Instalación de OpenCV.}]
pip3 install opencv-python-headless
\end{lstlisting}

\begin{figure}[h!]
    \centering
    \includegraphics[width=0.8\textwidth]{images/1_3actualizacion.png}
    \caption{Instalación de OpenCV.}
    \label{fig:opencv}
\end{figure}

Al estar trabajando con Windows, y ya tenía instalada la librería OpenCV de asignaturas previas, tenía algunos errores por lo que volví a la configuraciónn en la que no se usa ``headless''.


He accedido al repositorio oficial de código de OpenCV en GitHub (\textit{opencv/data/haarcascades}) para descargar de forma local los clasificadores necesarios. En la Figura \ref{fig:img1} se puede observar el contenido del repositorio. Para el ejercicio con reconocimiento de caras basta con descargar el fichero subrayado en naranja sobre rostros en posición frontal: \texttt{haarcascade\_frontalface\_alt.xml}. Para reconocer caras en videos hacen falta dos: \texttt{haarcascade\_eye\_tree\_eyeglasses.xml} y \texttt{haarcascade\_frontalface\_default.xml}.


\begin{figure}[ht]
    \centering
    \includegraphics[width=0.8\linewidth]{images/1_1descarga.png}
    \caption{Repositorio OpenCV}
    \label{fig:img1}
\end{figure}

\subsection{Ejemplo de caras}
El código fuente proporcionado por el profesor se encuentra en el Anexo \ref{anexo:ejemplo}. El material necesario (imagen y fichero \textit{.xml}) están en la siguiente dirección del repositorio: \url{https://github.com/ruiz314/PDIH/blob/main/P6/ejemplos/caras/ejemplo-fotos-1.py}.


La imagen inicial es la que se muestra en la Figura \ref{fig:img_ini}. Muestra dos rostros de dos hombres en un primer plano. Al fondo de la imagen hay un cuadro con un rostro de otro hombre. La imagen está en blanco y negro.

\begin{figure}[ht]
    \centering
    \includegraphics[width=0.7\linewidth]{images/imagen_input.jpg}
    \caption{Imagen inicial}
    \label{fig:img_ini}
\end{figure}

Para ejecutar este ejemplo en el que detecta las caras de una imagen basta con escribir \textit{python} y el nombre del fichero:

\begin{lstlisting}[style=terminalstyle]
python ejemplo-fotos-1.py
\end{lstlisting}

En la Figura \ref{fig:img4} se puede observar el resultado de haber ejecutado el ejemplo proporcionado por el profesor. Solo se detecta una cara, la de la derecha (Max Planck). En un recuadro amarillo que he dibujado en mi terminal se puede ver el contenido de la carpeta, y subrayado por una línea amarilla viene el comando para ejecutar dicho programa.


\begin{figure}[!ht]
    \centering
    \includegraphics[width=0.8\linewidth]{images/1_4ejecucion.png}
    \caption{Ejecución de ejemplo}
    \label{fig:img4}
\end{figure}

En la Figura \ref{fig:img5} se pueden ver dos recuadros rojos, es decir, el programa ha detectado dos caras en la imagen. La cara del cuadro no ha sido reconocida por el programa. Si reducimos el parámetro \texttt{scaleFactor} y le damos valor más cercano a $1.0$ (p.ej. $1.05$ en mi caso) veremos que sí identifica esa otra cara del cuadro. Subrayado en amarillo se encuentra el parámetro \texttt{scaleFactor}, para ver su valor en esa ejecución.

\begin{figure}[H]
    \centering
    \includegraphics[width=0.8\linewidth]{images/1_5ejec2.png}
    \caption{Ejecución de ejemplo con \texttt{scaleFactor = 1.05}}
    \label{fig:img5}
\end{figure}


%----------------------------------------------------------------
%% Sección IV
\newpage
\section{Ejercicios propuestos}\label{sec:Ejercicios}
% Ejercicios
\subsection{Ejercicio 1}\label{subsec:ejer1}
Para tener más imágenes con las que trabajar y probar el programa he buscado imágenes libres en la página \textit{Pixabay}: \url{https://pixabay.com/es/photos/search/fotos%20de%20personas/}. A continuación presento las imágenes con las que voy a trabajar:
\begin{itemize}
    \item Figura \ref{fig:imagen1} - Dos personas en un parque. La persona de la izquierda lleva una gorra, tiene barba y está mirando de frente. La persona de la derecha tiene el pelo corto y casi está mirando de frente a la cámara.
    \item Figura \ref{fig:imagen2} - Aparecen ocho personas, todas mirando a cámara y sonriendo.
\end{itemize}

\begin{figure}[ht]
    \centering
    \includegraphics[width=0.8\linewidth]{images/ejer1/imagen1.jpg}
    \caption{Imagen 1}
    \label{fig:imagen1}
\end{figure}

\begin{figure}[ht]
    \centering
    
\includegraphics[width=0.8\linewidth]{images/ejer1/imagen2.jpg}
    \caption{Imagen 2}
    \label{fig:imagen2}
\end{figure}

Para esta tarea se ha empleado el clasificador frontal avanzado \texttt{haarcascade\_frontalface\_alt.xml}. También importo la librería \texttt{sys} para poder incluir la imagen con la que se trabaja como parámetro en la orden de ejecución. Esto facilita probar con varias imágenes sin tener que tocar el fichero de código fuente. Además, uso la librería \texttt{os} para guardar las imágenes resultantes (con caras recuadradas) en la carpeta de trabajo.

\begin{lstlisting}[style=pythonstyle, caption={Librerías para el Ejercicio 1}, label={lst:lib1}]
import cv2
import sys
import os
\end{lstlisting}

Como la imagen con la que se va a trabajar se pasa por parámetro hay que comprobar que se haya pasado, que la ruta exista y que no esté corrupta:
\begin{lstlisting}[style=pythonstyle, caption={Comprobación de parametro}, label={lst:param1}]
# Validación del número de argumentos en la línea de comandos
if len(sys.argv) < 2:
    print("[ERROR] Uso correcto: python ejercicio1.py <nombre_imagen.jpg>")
    sys.exit(1)

# Captura del parámetro del archivo de imagen
nombre_imagen = sys.argv[1]

#Cargamos nuestro modelo classificador de Haar:
modelo = cv2.CascadeClassifier('haarcascade_frontalface_alt.xml')
  
#Cargamos la imagen
img = cv2.imread(nombre_imagen)

# Control de excepciones ante rutas inexistentes o corruptas
if img is None:
    print(f"[ERROR] No se pudo cargar la imagen '{nombre_imagen}'")
    sys.exit(1)    
\end{lstlisting}

A continuación, el programa sigue como el del ejemplo, haciendo el procesamiento de escala de grises, buscando rostros y dibujando los rectángulos sobre la imagen original. Una vez hecho eso, se guarda la imagen en disco separando el nombre de la imagen de la extensión:
\begin{lstlisting}[style=pythonstyle, caption={Guardar imagen}, label={lst:guardarimg1}]
# Separamos el nombre del archivo de su extensión (ej: 'foto' y '.jpg')
nombre_base, extension = os.path.splitext(nombre_imagen)

# Creamos el nuevo nombre para el archivo de salida
nombre_salida = f"resultado_{nombre_base}{extension}"

# Guardamos la imagen modificada en la carpeta de trabajo
cv2.imwrite(nombre_salida, img)
print(f"[INFO] Imagen resultante guardada con éxito como: {nombre_salida}")
\end{lstlisting}

El resultado usando \texttt{scaleFactor=1.35, minNeighbors=6} en la Figura \ref{fig:imagen1} es la Figura \ref{fig:1res_img1}. El motivo de usar \texttt{scaleFactor=1.35} es porque con el valor $1.5$ no se reconoce ningún rostro. El rostro del hombre no lo llega a reconocer aunque se disminuya el valor de \texttt{scaleFactor}.\\

\begin{figure}[ht]
    \centering
    \includegraphics[width=0.8\linewidth]{images/ejer1/resultado_imagen1.jpg}
    \caption{Resultado de \textit{imagen1} con \texttt{scaleFactor=1.35, minNeighbors=6}}
    \label{fig:1res_img1}
\end{figure}

El resultado usando \texttt{scaleFactor=1.5, minNeighbors=6} en la Figura \ref{fig:imagen2} es la Figura \ref{fig:1res_img2}. Se reconocen 5 rostros, de 8 que hay en la imagen.\\

\begin{figure}[ht]
    \centering
    \includegraphics[width=0.8\linewidth]{images/ejer1/resultado_imagen2.jpg}
    \caption{Resultado de \textit{imagen2} con \texttt{scaleFactor=1.5, minNeighbors=6}}
    \label{fig:1res_img2}
\end{figure}

\begin{figure}[ht]
    \centering
    \includegraphics[width=0.8\linewidth]{images/ejer1/resultado_imagen2_2.jpg}
    \caption{Resultado de \textit{imagen2} con \texttt{scaleFactor=1.35, minNeighbors=6}}
    \label{fig:1res_img2_2}
\end{figure}
Si disminuimos el valor de \texttt{scaleFactor} a $1.35$ el programa reconocerá todas las caras, como se muestra en la Figura \ref{fig:1res_img2_2}.\\


%----------------------------------------------------------------

\newpage
\section{Anexo}\label{anexo:an}

\subsection{Fichero \textit{ejemplo-fotos-1.py}} \label{anexo:ejemplo} 
Ejemplo proporcionado por el profesor. En el repositorio de \textit{GitHub} está en la siguiente ruta: \url{https://github.com/ruiz314/PDIH/blob/main/P6/ejemplos/caras/ejemplo-fotos-1.py}.
\lstinputlisting[style=pythonstyle, caption={Código completo del ejemplo1}, label={lst:ejemplo1}]{src/ejemplo-fotos-1.py}

% --- Ejercicio 1 
\subsection{Fichero \textit{ejercicio1.py}} \label{anexo:ejer1} 
\lstinputlisting[style=pythonstyle, caption={Código completo del ejercicio1}, label={lst:ejerc1}]{src/ejercicio1.py}

\subsection{Fichero \textit{ejemplo2.py}} \label{anexo:ejemplo2} 
Ejemplo proporcionado por el profesor. En el repositorio de \textit{GitHub} está en la siguiente ruta: \url{https://github.com/ruiz314/PDIH/blob/main/P6/ejemplos/en-videos/ejemplo2.py}.
\lstinputlisting[style=pythonstyle, caption={Código completo del ejemplo2}, label={lst:ejemplo2}]{src/ejemplo2.py}

\subsection{Fichero \textit{ejemplo3.py}} \label{anexo:ejemplo3} 
Ejemplo proporcionado por el profesor. En el repositorio de \textit{GitHub} está en la siguiente ruta: \url{https://github.com/ruiz314/PDIH/blob/main/P6/ejemplos/en-videos/ejemplo3.py}.
\lstinputlisting[style=pythonstyle, caption={Código completo del ejemplo3}, label={lst:ejemplo3}]{src/ejemplo3.py}

\end{document}
